\documentclass[11pt]{letter}

\usepackage[T1]{fontenc}
\usepackage[utf8]{inputenc}
\usepackage{lmodern}
\usepackage[margin=1in]{geometry}
\usepackage[hidelinks]{hyperref}
\usepackage{microtype}

\setlength{\parskip}{0.65em}
\setlength{\parindent}{0pt}

\signature{Christopher Robinson Tompkins III\\Corresponding author}
\address{%
    Christopher Robinson Tompkins III\\
    PhD Student Systems Modeling and Analysis, MS Mathematical Logic (academic descendant of Alfred Tarski)\\
    Richmond, Virginia, USA\\
    \href{mailto:chtompki@apache.org}{chtompki@apache.org}\\
    \href{mailto:chtompki@vt.edu}{chtompki@vt.edu}\\
    \href{mailto:tompkinscr@vcu.edu}{tompkinscr@vcu.edu}%
}

\begin{document}

    \begin{letter}{Editors\\Science}

        \date{August 7, 2026}
        \opening{Dear Editors,}

        Please consider my manuscript, ``On the Basis of Mathematics,'' for
        publication in \textit{Science}.

        The manuscript proposes ``writability''---the capacity of an idea or
        statement to be instantiated in communicable symbols---as a fundamental
        bound on mathematical language. Beginning with self-reference and the liar
        paradox, it develops a minimal representational framework from which
        reference, truth values, logical connectives, and a language of discourse can
        be constructed. It then examines a provocative consequence: if every
        mathematically usable expression must be writable, the conventional
        separation between object-language and meta-language may require
        reconsideration. From this perspective, the manuscript argues that
        mathematical completeness and incompleteness can be understood as
        simultaneously arising within a self-referential language.

        This work addresses foundational questions extending across mathematical
        logic, the philosophy of mathematics, theories of computation, and the
        scientific study of language and representation. Its principal contribution
        is to identify writability as an explicit constraint on formal inquiry and to
        use that constraint to offer a unified interpretation of self-reference,
        G\"odelian completeness and incompleteness, and Hilbert's second problem.
        Because the argument concerns the representational foundations shared by
        mathematics and the sciences, I believe it may interest \textit{Science}'s
        broad interdisciplinary readership.

        I am the sole author and corresponding author:

        \begin{quote}
            Christopher Robinson Tompkins III\\
            PhD Student Systems Modeling and Analysis, MS Mathematical Logic (academic descendant of Alfred Tarski)\\
            Richmond, Virginia, USA\\
            Email: \href{mailto:chtompki@vt.edu}{chtompki@vt.edu}\\
            Additional emails: \href{mailto:chtompki@apache.org}{chtompki@apache.org};
            \href{mailto:tompkinscr@vcu.edu}{tompkinscr@vcu.edu}
        \end{quote}

        This manuscript has not been published previously and
        is not under consideration by another journal. No closely related manuscript
        is under consideration or in press elsewhere.

        I have no competing interests to declare.

        For transparency, ChatGPT was used in developing certain examples and the
        idea of employing photographs to illustrate self-reference, as acknowledged
        in the manuscript. ChatGPT was also used as a writing aid in preparing this
        cover letter. I take full responsibility for the manuscript's arguments,
        claims, presentation, and submitted text.

        Thank you for considering this work. I would welcome the opportunity to have
        it evaluated by experts in mathematical logic, foundations of mathematics,
        and the philosophy of formal systems.

        \closing{Sincerely,}

    \end{letter}
\end{document}
